\documentclass[11pt,a4paper]{article}

% ---------- Packages ----------
\usepackage[utf8]{inputenc}
\usepackage[T1]{fontenc}
\usepackage{lmodern}
\usepackage[margin=1in]{geometry}
\usepackage{booktabs}
\usepackage{array}
\usepackage{xcolor}
\usepackage{titlesec}
\usepackage{enumitem}
\usepackage{caption}
\usepackage{fancyhdr}
\usepackage{listings}
\usepackage[hidelinks]{hyperref}

% ---------- Colors ----------
\definecolor{primary}{HTML}{1D4ED8}
\definecolor{slate}{HTML}{334155}
\definecolor{codebg}{HTML}{F1F5F9}

\hypersetup{
  colorlinks=true,
  linkcolor=primary,
  urlcolor=primary,
  citecolor=primary
}

% ---------- Headers / footers ----------
\pagestyle{fancy}
\fancyhf{}
\renewcommand{\headrulewidth}{0.4pt}
\lhead{\small\textcolor{slate}{Pearls AQI Predictor}}
\rhead{\small\textcolor{slate}{Project Report}}
\cfoot{\thepage}

% ---------- Section styling ----------
\titleformat{\section}{\Large\bfseries\color{primary}}{\thesection}{0.6em}{}
\titleformat{\subsection}{\large\bfseries\color{slate}}{\thesubsection}{0.6em}{}

% ---------- Code listing style ----------
\lstset{
  backgroundcolor=\color{codebg},
  basicstyle=\ttfamily\small,
  breaklines=true,
  frame=single,
  rulecolor=\color{codebg},
  framesep=6pt,
  xleftmargin=6pt,
  showstringspaces=false
}

\newcommand{\rsq}{$R^2$}

% ---------- Title ----------
\title{\vspace{-1cm}\textbf{\color{primary}\Huge Pearls AQI Predictor}\\[0.4cm]
{\Large A 100\% Serverless End-to-End Machine Learning System\\
for Forecasting Air Quality in Karachi}}
\author{Submitted for the 10Pearls AQI Predictor Assignment}
\date{\today}

\begin{document}

\maketitle
\thispagestyle{empty}

\begin{abstract}
\noindent
This report documents the design and implementation of an end-to-end machine
learning system that forecasts the Air Quality Index (AQI) for Karachi,
Pakistan over a three-day horizon (+24h, +48h, +72h). The system is fully
serverless and runs entirely on free tiers: raw weather and pollutant data are
fetched from the Open-Meteo API, engineered into 29 features, and persisted in a
Supabase feature store; five model families spanning statistical to deep-learning
approaches are trained and benchmarked, with champions promoted to a git-based
model registry; GitHub Actions automate the feature pipeline hourly and the
training pipeline daily; and an interactive Streamlit dashboard loads features
and models from the store to display live and forecasted AQI with SHAP-based
explanations and hazardous-level alerts. Ridge regression emerged as the
champion at all three horizons, achieving an RMSE of 21.88 at +24h on a
chronological hold-out set.
\end{abstract}

\vspace{0.5cm}
\begin{center}
\begin{tabular}{rl}
\textbf{Live dashboard:} & \url{https://10pearls-karachi-aqi-predictor-v8ygbuwud8du7m3mun22cr.streamlit.app}\\
\textbf{Source code:} & \url{https://github.com/Zee9710/10Pearls-Karachi-AQI-Predictor}\\
\textbf{Location:} & Karachi (24.8607\textdegree{} N, 67.0011\textdegree{} E), \texttt{Asia/Karachi}\\
\end{tabular}
\end{center}

\newpage
\tableofcontents
\newpage

% =====================================================================
\section{Introduction}
% =====================================================================
Air pollution is a persistent public-health concern in Karachi, one of the
world's most populous cities. The objective of this project is to predict the
Air Quality Index (AQI) for the next three days using a fully automated,
serverless machine learning pipeline. The system must fetch raw data, engineer
features, store them in a feature store, train and evaluate a range of models,
register the best performers, automate the entire workflow with CI/CD, and
surface predictions through an interactive dashboard.

The solution adopts the canonical MLOps pattern built around a
\textbf{feature store} and a \textbf{model registry}. Each stage is decoupled:
the feature pipeline only writes features, the training pipeline only reads
features and writes models, and the dashboard only reads. This separation makes
every component independently testable and replaceable, and allows the feature
and training pipelines to run on independent schedules.

% =====================================================================
\section{System Architecture}
% =====================================================================
Figure-style overview of the data flow:

\begin{center}
\begin{lstlisting}[language=,frame=single]
Open-Meteo API --> Feature Pipeline --> Feature Store (Supabase)
   (raw data)        feature_pipeline.py     aqi_features table
                                                    |
                                  +-----------------+------------------+
                                  |                                    |
                          Training Pipeline                    Streamlit Dashboard
                          training_pipeline.py                  streamlit_app.py
                                  |                                    ^
                                  v                                    |
                          Model Registry  ----------------------------+
                          (models/ in git)

GitHub Actions:  feature pipeline hourly  *  training pipeline daily
\end{lstlisting}
\end{center}

\subsection{Serverless Stack}
Every component runs on a free tier with no server to maintain.

\begin{center}
\begin{tabular}{@{}p{3.2cm}p{4.2cm}p{6.0cm}@{}}
\toprule
\textbf{Concern} & \textbf{Choice} & \textbf{Rationale} \\
\midrule
Raw data & Open-Meteo API & Free, no API key; provides air-quality, weather, and a historical archive endpoint \\
Feature store & Supabase (Postgres + PostgREST) & Free tier, SQL-queryable, accessed over plain HTTPS \\
Model registry & \texttt{models/} in git & Zero-cost, versioned, deploys with the app \\
Orchestration & GitHub Actions & Free cron scheduling, no server \\
Dashboard & Streamlit Community Cloud & Free hosting with a direct public URL \\
\bottomrule
\end{tabular}
\end{center}

% =====================================================================
\section{Feature Pipeline}
% =====================================================================
The feature pipeline (\texttt{src/feature\_pipeline.py}) performs four steps.

\subsection{Fetch Raw Data}
Hourly pollutant data (PM2.5, PM10, ozone, NO\textsubscript{2},
SO\textsubscript{2}, CO) is pulled from Open-Meteo's air-quality endpoint, and
weather data (temperature, relative humidity, wind speed, precipitation,
surface pressure) from the forecast/archive endpoints. A retry helper backs off
on HTTP 429 and 5xx responses and honors the \texttt{Retry-After} header, which
is important on the shared IP addresses used by cloud hosts.

\subsection{Compute the AQI Target}
The US-EPA AQI (0--500) is computed from raw pollutant concentrations in
\texttt{src/aqi.py} using the standard breakpoint piecewise-linear formula,
taking the maximum sub-index across pollutants.

\subsection{Engineer Features}
Twenty-nine model inputs are derived:

\begin{itemize}[leftmargin=1.4em]
  \item \textbf{Cyclical time:} \texttt{hour\_sin/cos}, \texttt{dow\_sin/cos},
        \texttt{month\_sin/cos}, \texttt{is\_weekend}. Sine/cosine encoding keeps
        23:00 and 00:00 adjacent in feature space.
  \item \textbf{Lags:} \texttt{aqi\_lag\_1h}, \texttt{aqi\_lag\_24h},
        \texttt{aqi\_lag\_48h}, \texttt{aqi\_lag\_72h}.
  \item \textbf{Rolling statistics:} 24h and 72h rolling mean, standard
        deviation, and maximum (shifted to prevent leakage).
  \item \textbf{Change rate:} \texttt{aqi\_change\_24h}, the day-over-day AQI
        delta required by the assignment.
  \item \textbf{Interactions:} \texttt{temp\_humidity}, \texttt{wind\_pollution}.
  \item \textbf{Raw signals:} all six pollutants and five weather variables.
\end{itemize}

\subsection{Store}
Features are upserted into the Supabase \texttt{aqi\_features} table, keyed on
\texttt{datetime}, so re-runs are idempotent and never create duplicate rows.

% =====================================================================
\section{Backfill}
% =====================================================================
To generate training data, the same feature script is run over a one-year range
of past dates (\texttt{src/backfill.py}). This populated the feature store with
approximately one year of hourly rows (around 8,900 records). The backfill uses
the Open-Meteo \emph{archive} endpoint for dates older than seven days and the
\emph{forecast} endpoint for recent dates.

% =====================================================================
\section{Training Pipeline}
% =====================================================================
The training pipeline (\texttt{src/training\_pipeline.py}) executes three steps.

\subsection{Fetch from the Feature Store}
The full feature history is read from Supabase. The training stage never touches
the live raw API, satisfying the requirement that models train on stored
features.

\subsection{Train and Evaluate}
For each horizon $h \in \{24, 48, 72\}$, the target is the AQI shifted forward
by $h$ hours, giving \emph{direct multi-horizon} forecasting. The data is split
80/20 \emph{chronologically} (no shuffling) to respect time ordering and prevent
look-ahead leakage. Five model families are trained:

\begin{center}
\begin{tabular}{@{}ll@{}}
\toprule
\textbf{Family} & \textbf{Type} \\
\midrule
Ridge Regression (\texttt{RidgeCV}) & Statistical / linear \\
Random Forest & Bagged decision trees \\
XGBoost & Gradient boosting \\
LightGBM & Gradient boosting \\
LSTM (TensorFlow/Keras) & Deep learning (sequence model) \\
\bottomrule
\end{tabular}
\end{center}

This spans the spectrum from statistical modelling to deep learning. Each model
is scored on RMSE, MAE, and \rsq{}; the lowest-RMSE model per horizon is
promoted to champion.

\subsection{Register}
The champion model, its scaler, and its metrics are written to
\texttt{models/aqi\_best\_\{24,48,72\}h/} and committed to git by the daily CI
job.

% =====================================================================
\section{Results}
% =====================================================================
AQI is on the 0--500 EPA scale; all metrics are computed on the chronological
20\% hold-out tail.

\subsection{Champion per Horizon}
\begin{center}
\begin{tabular}{@{}lcccc@{}}
\toprule
\textbf{Horizon} & \textbf{Champion} & \textbf{RMSE} & \textbf{MAE} & \textbf{\rsq{}} \\
\midrule
+24h & Ridge & 21.88 & 15.53 & 0.308 \\
+48h & Ridge & 24.59 & 18.07 & 0.128 \\
+72h & Ridge & 25.47 & 18.99 & 0.065 \\
\bottomrule
\end{tabular}
\end{center}

\subsection{Full Model Comparison (RMSE)}
\begin{center}
\begin{tabular}{@{}lccc@{}}
\toprule
\textbf{Model} & \textbf{+24h} & \textbf{+48h} & \textbf{+72h} \\
\midrule
\textbf{Ridge} & \textbf{21.88} & \textbf{24.59} & \textbf{25.47} \\
Random Forest  & 22.95 & 28.86 & 28.05 \\
XGBoost        & 23.74 & 28.30 & 27.10 \\
LightGBM       & 24.24 & 28.64 & 27.49 \\
LSTM           & 29.81 & 29.90 & 29.60 \\
\bottomrule
\end{tabular}
\end{center}

\subsection{Discussion}
\begin{itemize}[leftmargin=1.4em]
  \item \textbf{Ridge wins at every horizon.} With roughly 8,900 hourly rows and
        strongly autocorrelated lag features, the linear model generalizes
        better than the tree ensembles and the LSTM, which overfit (negative
        \rsq{} on the held-out tail at the longer horizons).
  \item \textbf{Accuracy decays with horizon}, as expected: +24h is the most
        predictable (\rsq{} $\approx$ 0.31), while +72h approaches the
        persistence baseline (\rsq{} $\approx$ 0.07).
  \item \textbf{The LSTM is data-starved.} Deep sequence models typically need
        far more data than is currently available; the LSTM is retained to
        demonstrate the full statistical-to-deep-learning spectrum and should be
        re-benchmarked as the feature store grows hourly.
\end{itemize}

% =====================================================================
\section{CI/CD Automation}
% =====================================================================
Two GitHub Actions workflows automate the system end to end.

\begin{center}
\begin{tabular}{@{}p{3.6cm}p{3.2cm}p{6.4cm}@{}}
\toprule
\textbf{Workflow} & \textbf{Schedule} & \textbf{Action} \\
\midrule
\texttt{hourly-feature.yml} & \texttt{0 * * * *} (hourly) & Fetch the latest window, compute features, upsert to the store \\
\texttt{daily-train.yml} & \texttt{0 3 * * *} (daily) & Read the store, retrain all models, commit updated champions \\
\bottomrule
\end{tabular}
\end{center}

Both authenticate to Supabase via repository secrets
(\texttt{SUPABASE\_URL}, \texttt{SUPABASE\_KEY}) and require no server. The
training job installs \texttt{requirements-train.txt} (adding TensorFlow on top
of the base requirements) and pushes refreshed models back to the repository.

% =====================================================================
\section{Web Application}
% =====================================================================
The dashboard (\texttt{streamlit\_app.py}) is deployed on Streamlit Community
Cloud and styled after the IQAir interface.

\subsection{Loads Model and Features from the Feature Store}
On load, the app reads the most recent feature rows directly from Supabase and
loads the champion models from the registry. If credentials are absent it falls
back to a live Open-Meteo fetch so the demo still works; the active source is
shown in the interface.

\subsection{Displays Predictions}
The dashboard presents:
\begin{itemize}[leftmargin=1.4em]
  \item Current AQI with a category ring, health advisory, and WHO PM2.5
        multiplier.
  \item A live status pill showing the data source and last update time.
  \item Six pollutant concentration cards.
  \item Three-day forecast cards (+24h / +48h / +72h) with color-coded
        categories.
  \item A historical AQI trend chart (7 / 14 / 30 day windows).
  \item A model-performance section comparing RMSE across all five models per
        horizon, with the champion highlighted.
  \item A SHAP feature-importance section explaining which inputs drive each
        forecast.
\end{itemize}

\subsection{Hazardous-AQI Alerts}
When a forecast crosses the alert threshold (AQI 150, ``Unhealthy''), the
dashboard surfaces a prominent warning.

% =====================================================================
\section{Exploratory Data Analysis and Guidelines}
% =====================================================================
The notebook \texttt{notebooks/eda.ipynb} examines pollutant distributions,
feature correlations, and time-of-day and seasonal trends in the AQI signal,
informing the choice of lag and rolling features. The table below maps each
assignment guideline to its implementation.

\begin{center}
\begin{tabular}{@{}p{6.0cm}p{8.0cm}@{}}
\toprule
\textbf{Guideline} & \textbf{Implementation} \\
\midrule
EDA to identify trends & \texttt{notebooks/eda.ipynb} \\
Variety of models (statistical $\rightarrow$ deep learning) & Ridge, Random Forest, XGBoost, LightGBM, LSTM \\
SHAP / LIME explainability & SHAP in the dashboard; LIME in the optional Flask API \\
Alerts for hazardous AQI & Threshold alert at AQI 150 \\
\bottomrule
\end{tabular}
\end{center}

% =====================================================================
\section{Deliverables Checklist}
% =====================================================================
\begin{enumerate}[leftmargin=1.6em]
  \item \textbf{End-to-end AQI prediction system} --- complete
        (data $\rightarrow$ features $\rightarrow$ models $\rightarrow$ predictions).
  \item \textbf{Scalable, automated pipeline} --- complete (hourly and daily
        GitHub Actions; idempotent feature store).
  \item \textbf{Interactive dashboard} --- complete (live and forecasted AQI on
        Streamlit Cloud).
  \item \textbf{Detailed report} --- this document.
\end{enumerate}

% =====================================================================
\section{Limitations and Future Work}
% =====================================================================
\begin{itemize}[leftmargin=1.4em]
  \item \textbf{Forecast skill at long horizons.} +72h \rsq{} is low; richer
        exogenous signals (wind direction, regional transport, traffic) could
        improve it.
  \item \textbf{LSTM needs more data.} As the hourly pipeline accumulates
        history, the deep model should be re-benchmarked.
  \item \textbf{Single location.} The pipeline is parameterized by latitude and
        longitude in \texttt{src/config.py} and could be extended to multiple
        cities.
  \item \textbf{Git-based model registry.} Adequate at this scale; a managed
        registry such as MLflow would be the next step for model lineage in a
        larger team.
\end{itemize}

\vspace{0.5cm}
\hrule
\vspace{0.3cm}
\noindent\small\textit{Reproduce locally:}
\begin{lstlisting}[language=bash]
pip install -r requirements-train.txt
python -m src.backfill            # populate the feature store
python -m src.training_pipeline   # train + register champions
streamlit run streamlit_app.py    # launch the dashboard
\end{lstlisting}

\end{document}
